\section{A common coefficient algebra for the nearby original normal cones}

The coefficient algebra here consists of actual functions holomorphic near every
point of a fixed open base set. It is distinct from the ring of germs at the
origin: passage to each nearby germ is defined on the actual functions. This
supplies a common algebraic presentation of the previously identified nearby
normal cones. All-degree generic freeness for this presentation is established in the
subsequent base-change and nearby-freeness sections.
This coordinate normal-module calculation is an additional commutative-algebra
step in the resolution route of W\l odarczyk's marked-ideal construction.
The common function algebra, whole-kernel shrinking, and uniform polynomial
denominator below are the explicit analytic adaptation used here; they do not
assert the geometric singular-locus-centre or higher Hilbert--Samuel driver.


\begin{definition}[The algebra of holomorphic coefficient functions]
\label{hir:relative-normal-coefficient-function-algebra}
\source{wlodarczyk05-hironaka}
\lean{Hironaka.HironakaRelativeNormalCoefficientFunctionAlgebra.coefficientFunctionAlgebra}
\leanok
For a base set $B\subset\mathbb C^d$, let $A_B$ be the complex subalgebra of
functions $\mathbb C^d\to\mathbb C$ that are holomorphic in a neighbourhood of
each point of $B$. Its addition and multiplication are pointwise operations.
\end{definition}
\begin{proof}
Constant functions are holomorphic. The sum, product and negative of functions
holomorphic near a point are holomorphic near that point, by their complex
derivative formulas on the intersection of the two neighbourhoods. These
closures give the asserted complex subalgebra.
\end{proof}

\begin{lemma}[Nearby germ maps on the common coefficient algebra]
\label{hir:relative-normal-coefficient-function-germ-hom}
\source{wlodarczyk05-hironaka}
\lean{Hironaka.HironakaRelativeNormalCoefficientFunctionAlgebra.coefficientGermHom}
\leanok
\uses{hir:relative-normal-coefficient-function-algebra,hir:relative-normal-nearby-actual-function-coefficient-germs}
For every $a\in B$ there is a complex algebra homomorphism
$\rho_a:A_B\to\mathcal O_{\mathbb C^d,0}$ sending an actual function $f$ to
the whole germ of $z\mapsto f(a+z)$.
\end{lemma}
\begin{proof}
Holomorphicity of $f$ near $a$ and the complex chain rule make its translated
function holomorphic near zero. Translation commutes pointwise with constants,
sums and products, and taking a whole germ preserves these operations.
Consequently this formula defines the asserted algebra homomorphism.
\end{proof}

\begin{lemma}[Whole polynomial lifts over the actual coefficient algebra]
\label{hir:relative-normal-coefficient-function-polynomial-lift}
\source{wlodarczyk05-hironaka}
\lean{Hironaka.HironakaRelativeNormalCoefficientFunctionAlgebra.coefficientPolynomialLift,Hironaka.HironakaRelativeNormalCoefficientFunctionAlgebra.coefficientPolynomialLift_coeff,Hironaka.HironakaRelativeNormalCoefficientFunctionAlgebra.coefficientPolynomialLift_homogeneous,Hironaka.HironakaRelativeNormalCoefficientFunctionAlgebra.coefficientPolynomialLift_nearby,Hironaka.HironakaRelativeNormalCoefficientFunctionAlgebra.coefficientPolynomialLift_origin}
\leanok
\uses{hir:relative-normal-coefficient-function-germ-hom,hir:relative-normal-analytic-coefficients-whole-polynomial-dictionary,hir:relative-normal-nearby-entire-polynomial-coefficient-germs}
Suppose the chosen representative of every coefficient of $P$ is holomorphic
near every point of $B$. The finite sum with those same coefficient functions
defines $\widetilde P\in A_B[Y_1,\ldots,Y_t]$. At every exponent its
coefficient is the actual chosen coefficient function of $P$. Thus its support
is contained in that of $P$, and its homogeneous degree is retained. For every
$a\in B$, mapping $\widetilde P$ by $\rho_a$ gives the entire nearby
coefficient polynomial $P_a$. If $0\in B$, mapping by $\rho_0$ gives $P$.
\end{lemma}
\begin{proof}
Use the sum over the finite support of $P$. Coefficient extraction gives the
asserted coefficient at a supported exponent. Outside that support the original
coefficient is zero and its chosen function is identically zero, so the same
formula holds there. Every nonzero coefficient of the lift therefore has an
original supported exponent, which proves preservation of homogeneous degree.
Extract coefficients after applying each algebra homomorphism: they are exactly
the translated whole germs used to define $P_a$. At zero the translation is the
identity and the chosen function has its prescribed original whole germ.
Polynomial equality follows at every exponent.
\end{proof}

\begin{theorem}[One actual algebra presents every entire nearby original initial ideal]
\label{hir:relative-normal-coefficient-function-original-initial-family}
\source{wlodarczyk05-hironaka}
\lean{Hironaka.HironakaRelativeNormalCoefficientFunctionAlgebra.exists_original_coefficient_algebra_initial_family}
\leanok
\uses{hir:relative-normal-coefficient-function-polynomial-lift,hir:relative-normal-affine-dnc-original-nearby-cone,hir:relative-normal-graded-internal-finite-relations}
An original coherent ideal $I$ internally supplies one open base neighbourhood
$B$ of zero and lifts over $A_B$ of its selected finite original homogeneous
relation family. If $K_B$ is the ideal generated by those lifts, its image under
the polynomial map induced by $\rho_0$ is the whole original normal initial
ideal. At every $a\in B$, its image under the polynomial map induced by
$\rho_a$ is the whole normal initial ideal of the actual translated original
ideal. The generators over $A_B$ are homogeneous.
\end{theorem}
\begin{proof}
The entire nearby original normal-cone theorem supplies a single open $B$ on
which every selected coefficient function is holomorphic near every base point
and on which the nearby finite relation polynomials generate precisely the
whole original initial ideals. Form their lifts by the preceding lemma. The
image of an ideal generated by a family is the ideal generated by the images
of that family. At each $a$ those images are exactly the nearby polynomials,
so the existing entire nearby initial-ideal equality gives the desired
equality. At zero each image is its original relation polynomial, and the
selected original finite relations generate the whole original initial ideal.
Each selected relation has its original homogeneous degree; the preceding
lemma retains that degree over $A_B$.
\end{proof}

\begin{lemma}[Finite type coefficient algebra for finitely many polynomials]
\label{hir:relative-normal-finite-type-coefficient-algebra}
\source{wlodarczyk05-hironaka}
\lean{Hironaka.HironakaRelativeNormalCoefficientFunctionAlgebra.finitePolynomialCoefficientAlgebra,Hironaka.HironakaRelativeNormalCoefficientFunctionAlgebra.finitePolynomialCoefficientAlgebra_finiteType,Hironaka.HironakaRelativeNormalCoefficientFunctionAlgebra.finitePolynomialCoefficientAlgebra_isNoetherianRing,Hironaka.HironakaRelativeNormalCoefficientFunctionAlgebra.finiteCoefficientPolynomialLift_map,Hironaka.HironakaRelativeNormalCoefficientFunctionAlgebra.finiteCoefficientPolynomialLift_homogeneous}
\leanok
\uses{hir:relative-normal-coefficient-function-original-initial-family}
For finitely many polynomials over $A_B$, the complex subalgebra $C$ generated
by their finitely many supported coefficients is of finite type and
Noetherian. Each polynomial lifts to $C[Y_1,\ldots,Y_t]$ with exactly its
original image under coefficient inclusion, retaining homogeneous degree.
\end{lemma}
\begin{proof}
The union, over the finite family, of the coefficient sets indexed by each
finite support is finite. Its complex algebraic span is therefore a finite
type algebra and a quotient of a polynomial algebra in finitely many variables
over the Noetherian field $\mathbb C$, hence Noetherian. Every supported
coefficient lies in this subalgebra by construction; unsupported coefficients
are zero and also lie there. Form the same finite sum over each support with
these coefficients in $C$. Inclusion recovers every original coefficient,
and injectivity of inclusion shows that the lifted support is exactly the
original support. Homogeneous degree is consequently unchanged.
\end{proof}

\begin{theorem}[Original nearby initial ideals over one finite type coefficient algebra]
\label{hir:relative-normal-original-finite-type-coefficient-initial-family}
\source{wlodarczyk05-hironaka}
\lean{Hironaka.HironakaRelativeNormalCoefficientFunctionAlgebra.exists_original_finiteType_coefficient_algebra_initial_family}
\leanok
\uses{hir:relative-normal-coefficient-function-original-initial-family,hir:relative-normal-finite-type-coefficient-algebra}
Only the original coherent ideal $I$ supplies an open $B$ containing zero,
a finite type complex subalgebra $C\subset A_B$, and a finite homogeneous
relation family over $C$. Its entire ideal image under the actual origin germ
map is the original normal initial ideal. Its image under the actual nearby
germ map at every $a\in B$ is the entire normal initial ideal of the
actual translated original ideal.
\end{theorem}
\begin{proof}
The original-input common function-algebra theorem gives $B$ and its finite
homogeneous relations over $A_B$. Take $C$ to be the complex subalgebra
generated by their finite supported coefficient set. The finite type lemma
lifts every relation to $C$, with its original image under inclusion and its
original homogeneous degree. Composing the origin and nearby maps with this
inclusion gives exactly the previous images of every generator. The ideal
map formula for generated ideals therefore retains each entire initial-ideal
equality, both at origin and at every point of $B$.
\end{proof}

\begin{lemma}[One neighbourhood kills the entire origin coefficient kernel]
\label{hir:relative-normal-coefficient-algebra-kernel-spread}
\source{wlodarczyk05-hironaka}
\lean{Hironaka.HironakaRelativeNormalCoefficientFunctionAlgebra.exists_common_coefficient_kernel_open}
\leanok
\uses{hir:relative-normal-finite-type-coefficient-algebra,hir:relative-normal-coefficient-function-germ-hom}
Let $B$ be open and contain zero, and let $C\subset A_B$ be a Noetherian
complex subalgebra. There is one open neighbourhood $W\subset B$ of zero
such that $\ker(\rho_0|_C)\subset\ker(\rho_a|_C)$ for every $a\in W$.
\end{lemma}
\begin{proof}
Noetherianity makes the origin kernel finitely generated. Each of its selected
generators is an actual function with zero whole germ at the origin, so it
vanishes on a neighbourhood of zero. Intersect these finitely many
neighbourhoods and $B$, and take an open neighbourhood $W$ of zero inside
that intersection. At every $a\in W$, each selected actual function is
identically zero near $a$. Its translated whole germ at zero is therefore
zero. The nearby homomorphism kills the selected generators; since its kernel
is an ideal, it kills their whole span, which is exactly the origin kernel.
\end{proof}

\begin{theorem}[A common domain coefficient algebra with genuine nearby homomorphisms]
\label{hir:relative-normal-coefficient-algebra-domain-nearby-homs}
\source{wlodarczyk05-hironaka}
\lean{Hironaka.HironakaRelativeNormalCoefficientFunctionAlgebra.coefficientOriginKernel,Hironaka.HironakaRelativeNormalCoefficientFunctionAlgebra.coefficientOriginAlgebra_isDomain,Hironaka.HironakaRelativeNormalCoefficientFunctionAlgebra.coefficientOriginAlgebra_finiteType,Hironaka.HironakaRelativeNormalCoefficientFunctionAlgebra.exists_common_coefficient_origin_quotient_homs}
\leanok
\uses{hir:relative-normal-coefficient-algebra-kernel-spread}
The quotient $C_0=C/\ker(\rho_0|_C)$ is a domain and, when $C$ is of
finite type over $\mathbb C$, remains of finite type. On one open neighbourhood
$W$ of zero every actual nearby homomorphism $\rho_a|_C$ factors as a complex
algebra homomorphism through $C_0$, with its original actual-function formula
preserved after composition with the quotient map.
\end{theorem}
\begin{proof}
The origin homomorphism takes values in the domain of holomorphic germs, so its
kernel is prime and the quotient is a domain. A quotient of a finite type
algebra is of finite type. The preceding lemma gives one $W$ on which this
entire origin kernel lies in every nearby kernel. The universal property of
the quotient therefore defines each asserted homomorphism, and its defining
composition identity recovers the actual translated whole germ map. No map
on arbitrary origin germs is required: only the finitely generated actual
coefficient algebra is factored through its internally determined kernel.
\end{proof}

\begin{theorem}[A single finite type domain presents all nearby original initial ideals]
\label{hir:relative-normal-original-common-domain-coefficient-initial-family}
\source{wlodarczyk05-hironaka}
\lean{Hironaka.HironakaRelativeNormalCoefficientFunctionAlgebra.exists_original_common_domain_coefficient_initial_family}
\leanok
\uses{hir:relative-normal-original-finite-type-coefficient-initial-family,hir:relative-normal-coefficient-algebra-domain-nearby-homs}
Only the original coherent ideal $I$ supplies a finite type domain
$C_0=C/\ker(\rho_0|_C)$, an open base neighbourhood $W$ of zero,
and one finite homogeneous relation family over $C_0$. There are genuine
complex algebra homomorphisms from $C_0$ to the translated base germ rings
at every $a\in W$. Their composition with the quotient map equals the
actual translated whole germ map on coefficient functions. At every
$a\in W$, the image of the ideal generated by this same relation family is
the entire normal initial ideal of the actual translated original ideal.
\end{theorem}
\begin{proof}
Apply the preceding original-input theorem to obtain $B,C$ and its finite
homogeneous relations. Finite type over $\mathbb C$ makes $C$ Noetherian.
The common-domain theorem produces one $W\subset B$ on which every
nearby actual coefficient map factors through $C_0$. Map each relation's
coefficients through $C\to C_0$. Coefficient maps preserve homogeneous
degree. At every $a\in W$, composition of polynomial coefficient maps
and the defining factorisation identity give exactly the original nearby
relation polynomial. Since ideal maps preserve generated ideals, the entire
nearby initial-ideal equality follows for the same finite family. The
origin kernel is prime and quotient finite type is preserved, giving the
claimed domain and finite type properties. This presentation alone does not
assert freeness of its degree pieces or analytic normal flatness.
\end{proof}
